\subsection{自由群胚}

定义 9. —— 设 C 为箭图。在 C 上构造的自由群胚（记作 \(\mathrm{Grp}(\mathrm{C})\) ）是群胚 \(\varpi_{\widetilde{C}}\) ——与 C 相伴的图 \(\widetilde{C}\) 的道路类群胚。

Grp(C) 的顶点集等于 C 的顶点集；Grp(C) 的轨道是 C 的连通分支。

设 C 为非空箭图。由第 174 页推论 4，与 C 相伴的图 \(\widetilde{\mathrm{C}}\) 是树，当且仅当在 C 上构造的自由群胚 \(\mathrm{Grp}(\mathrm{C})\) 是单传递的。

典范箭图态射 \(i\colon \mathrm{C}\to \widetilde{\mathrm{C}}\) 与 \(j\colon \widetilde{\mathrm{C}}\to \varpi_{\widetilde{\mathrm{C}}}\) 的合成是从 C 到 \(\mathrm{Grp}(\mathrm{C})\) 的箭图态射，我们称之为从 C 到 \(\mathrm{Grp}(\mathrm{C})\) 的典范态射。把它记作 \(\theta\) 。对 C 的每个顶点 \(a\) ，有 \(\theta (a) = a\) 。若 \(f\) 是 C 的箭，则 \(\theta (f)\) 是道路 \((o(f),(f,1),t(f))\) （属于 \(\widetilde{\mathrm{C}}\) ）的类。态射 \(\theta\) 是单射（II, 第 174 页, 推论 2），其像生成群胚 \(\mathrm{Grp}(\mathrm{C})\) 。

设 \(\varphi\colon\mathrm{C}\to\mathrm{C}^{\prime}\) 为箭图态射。以 \(\theta_{\mathrm{C}}\) （相应地，\(\theta_{\mathrm{C}^{\prime}}\) ）表示从 C 到 \(\mathrm{Grp}(\mathrm{C})\) （相应地，从 \(\mathrm{C}^{\prime}\) 到 \(\mathrm{Grp}(\mathrm{C}^{\prime})\) ）的典范态射。态射 \(\varpi(\widetilde{\varphi})\) 是唯一的群胚态射 \(\psi\) ，它从 \(\mathrm{Grp}(\mathrm{C})\) 到 \(\mathrm{Grp}(\mathrm{C}^{\prime})\) ，使得 \(\psi\circ\theta_{\mathrm{C}}=\theta_{\mathrm{C}^{\prime}}\circ\varphi\) 。

命题 6. —— 设 C 为箭图，G 为群胚，\(\varphi\) 为从 C 到 G 的箭图态射。存在唯一的从 Grp(C) 到 G 的群胚态射 \(\varphi'\) ，使得 \(\varphi' \circ \theta_{\mathrm{C}} = \varphi\) 。

设 \(\psi\) 为从 \(\widetilde{\mathbf{C}}\) 到 G 的箭图态射，满足 \(\psi(a) = \varphi(a)\) （对 C 的每个顶点 \(a\) ）以及 \(\psi(f, \varepsilon) = \varphi(f)^{\varepsilon}\) （对 C 的每条箭 \(f\) 及每个 \(\varepsilon \in \{-1, 1\}\) ）。我们有 \(\psi \circ i = \varphi\) 。由第 173 页命题 5，存在态射 \(\varphi'\) ，它从 \(\mathrm{Grp}(\mathrm{C}) = \varpi_{\widetilde{\mathrm{C}}}\) 到 G，使得 \(\varphi' \circ j = \psi\) 。于是有 \(\varphi' \circ \theta_{\mathrm{C}} = \varphi' \circ j \circ i = \psi \circ i = \varphi\) 。如同命题 5 的证明那样，\(\varphi'\) 的唯一性由 \(\theta_{\mathrm{C}}(\mathrm{C})\) 生成群胚 \(\mathrm{Grp}(\mathrm{C})\) 这一事实得出。

在命题 6 的假设下，有时称 \(\varphi'\) 为从 Grp(C) 到 G 的、开拓箭图态射 \(\varphi\) 的群胚态射。

例. —— 设 C 为只有一个顶点 s 的箭图，A 为其箭集。群胚 \(\mathrm{Grp}(\mathrm{C})\) 就是与在 A 上构造的自由群 \(\mathrm{F}(\mathrm{A})\) 相伴的群胚。事实上，它只有一个顶点 s，且由命题 6 立即推出，从 A 到 \(\mathrm{Grp}(\mathrm{C})_{s}\) 的典范映射——它由从 C 到 \(\mathrm{Grp}(\mathrm{C})\) 的典范态射导出——满足自由群的泛性质（A, I, 第 85 页, 命题 8）。

